\documentclass[11pt]{article}
\usepackage[a4paper,margin=28mm]{geometry}
\usepackage{amsmath,amssymb,amsthm,mathtools}
\usepackage{enumitem}
\usepackage[hidelinks]{hyperref}

\newtheorem{definition}{Definition}
\newtheorem{conjecture}{Conjecture}
\newtheorem{question}{Research Question}
\newtheorem{theorem}{Theorem}
\newtheorem{remark}{Remark}

\title{Polynomial-Time Approximation of the Star Height of a DFA}
\author{Yoav Danieli}
\date{}

\begin{document}
\maketitle

\section{Regular expressions and star height}

Fix a finite alphabet $\Sigma$.

\begin{definition}[Ordinary regular expression]
The regular expressions over $\Sigma$ are generated by
\[
 E ::= \emptyset\mid\varepsilon\mid a
 \mid E\cup E\mid E\cdot E\mid E^*,
 \qquad a\in\Sigma.
\]
Their semantics is
\[
\begin{aligned}
L(\emptyset)&=\emptyset,&
L(\varepsilon)&=\{\varepsilon\},&
L(a)&=\{a\},\\
L(E\cup F)&=L(E)\cup L(F),&
L(EF)&=\{uv:u\in L(E),\,v\in L(F)\},\\
L(E^*)&=\bigcup_{m\geq0}L(E)^m.
\end{aligned}
\]
\end{definition}

\begin{definition}[Star height]
The star height $h(E)$ of a regular expression is defined by
\[
\begin{aligned}
h(\emptyset)=h(\varepsilon)=h(a)&=0,\\
h(E\cup F)=h(EF)&=\max\{h(E),h(F)\},\\
h(E^*)&=1+h(E).
\end{aligned}
\]
For a regular language $L\subseteq\Sigma^*$, its star height is
\[
 h(L)=\min\{h(E):L(E)=L\}.
\]
\end{definition}

Since expressions without Kleene star denote finite languages,
\[
 h(L)=0 \quad\Longleftrightarrow\quad L\text{ is finite}.
\]
For a DFA, finiteness is decidable in polynomial time by checking whether a
directed cycle is reachable from the initial state and can reach an accepting
state.

\section{DFA input model}

\begin{definition}[Deterministic finite automaton]
A DFA is a tuple
\[
 A=(Q,\Sigma,\delta,q_0,F),
\]
where $Q$ is finite, $\delta:Q\times\Sigma\to Q$ is total,
$q_0\in Q$, and $F\subseteq Q$. Its language is
\[
 L(A)=\{w\in\Sigma^*:\delta(q_0,w)\in F\}.
\]
The input size parameter used below is $n=|Q|$. Unless stated otherwise, the
DFA is assumed to be trim and minimal.
\end{definition}

\section{Polynomial-size output: regular-expression circuits}

An explicit regular expression equivalent to an $n$-state DFA may require
exponential size. A polynomial-time algorithm therefore cannot always print
the fully unfolded expression. We use shared subexpressions.

\begin{definition}[Regular-expression circuit]
A regular-expression circuit is a finite directed acyclic graph with one
output node. Each node is labelled by
\[
 \emptyset,\quad\varepsilon,\quad a\in\Sigma,\quad
 \cup,\quad\cdot,\quad\text{or}\quad *,
\]
with the expected arity. Its semantics is obtained by evaluating the circuit
as a regular expression. Its size is its number of nodes.

The star height of the circuit is the maximum number of $*$-nodes on a
directed path from the output to a leaf. This equals the star height of its
unfolding.
\end{definition}

A \emph{constructive polynomial-time approximation algorithm} must output a
polynomial-size regular-expression circuit. An output-sensitive variant may
instead print an ordinary expression in time polynomial in the input and
output sizes.

\section{Kleene's dynamic-programming construction}

Let $Q=\{q_1,\ldots,q_n\}$. For $0\leq k\leq n$, let
$R_{ij}^{(k)}$ denote the words taking $q_i$ to $q_j$ whose intermediate
states belong to $\{q_1,\ldots,q_k\}$. The standard recurrence is
\[
 R_{ij}^{(k)}
 =
 R_{ij}^{(k-1)}
 \cup
 R_{ik}^{(k-1)}
 \bigl(R_{kk}^{(k-1)}\bigr)^*
 R_{kj}^{(k-1)}.
\]
The base expressions $R_{ij}^{(0)}$ are unions of transition labels, together
with $\varepsilon$ when $i=j$. The output is the union of
$R_{0f}^{(n)}$ over accepting states $f$.

With sharing, this gives a polynomial-size circuit. Syntactically,
\[
 h\bigl(R_{ij}^{(k)}\bigr)\leq k,
\]
up to an inessential indexing shift. Thus the construction gives height at
most $n$. For every infinite language, $h(L)\geq1$, so this is an
$O(n)$-factor constructive approximation.

The order of the states can affect the resulting expression substantially,
but merely choosing the best ordering of the input DFA does not necessarily
find the optimal star height of the language.

\section{Cycle rank and the automata characterization}

\begin{definition}[Cycle rank]
Let $G$ be a finite directed graph. Its cycle rank $r(G)$ is defined
recursively:
\begin{enumerate}[label=(\roman*)]
  \item if $G$ is acyclic, then $r(G)=0$;
  \item if $G$ is strongly connected and has at least one edge, then
  \[
    r(G)=1+\min_{v\in V(G)}r(G-v);
  \]
  \item otherwise, $r(G)$ is the maximum cycle rank of a strongly connected
        component of $G$.
\end{enumerate}
\end{definition}

For an automaton $B$, write $r(B)$ for the cycle rank of its transition graph.

\begin{theorem}[Eggan]
For every regular language $L$,
\[
 h(L)=\min\{r(B):B\text{ is an NFA accepting }L\}.
\]
\end{theorem}

In particular, for the input DFA $A$,
\[
 h(L(A))\leq r(A)\leq n.
\]
The first inequality may be strict because Eggan's minimum ranges over all
equivalent NFAs, not merely the input DFA or its state-elimination order.

\begin{definition}[Bideterministic automaton]
A trim DFA is bideterministic if it has one accepting state and reversing all
transitions, while exchanging initial and accepting states, yields a
deterministic automaton.
\end{definition}

For bideterministic languages, the minimal trim DFA satisfies
\[
 h(L(A))=r(A).
\]
Consequently, approximation of star height for this class is exactly
approximation of cycle rank.

\section{Approximation conventions}

A multiplicative approximation needs special treatment at optimum $0$.

\begin{definition}[Value approximation]
Let $\alpha:\mathbb N\to[1,\infty)$. A polynomial-time
$\alpha(n)$-approximation for DFA star height is an algorithm which, on input
an $n$-state minimal DFA $A$, outputs an integer $\widehat h(A)$ such that:
\begin{enumerate}[label=(\roman*)]
  \item if $h(L(A))=0$, then $\widehat h(A)=0$;
  \item if $h(L(A))\geq1$, then
  \[
    h(L(A))
    \leq \widehat h(A)
    \leq \alpha(n)\,h(L(A)).
  \]
\end{enumerate}
\end{definition}

\begin{definition}[Constructive approximation]
A constructive $\alpha(n)$-approximation outputs a polynomial-size
regular-expression circuit $C_A$ satisfying
\[
 L(C_A)=L(A),
\]
and
\[
\begin{cases}
h(C_A)=0,&h(L(A))=0,\\[2mm]
h(C_A)\leq\alpha(n)\,h(L(A)),&h(L(A))\geq1.
\end{cases}
\]
\end{definition}

Constructive approximation is stronger than value approximation because the
output circuit certifies the upper bound.

\begin{definition}[Best polynomial-time approximation ratio]
Let $\alpha_{\mathrm{val}}^*(n)$ be the pointwise infimum, up to asymptotic
constants, of the approximation ratios achieved by polynomial-time value
algorithms. Define $\alpha_{\mathrm{con}}^*(n)$ analogously for constructive
algorithms.

The central quantitative problem is to determine the asymptotic growth of
these functions.
\end{definition}

The Kleene construction gives
\[
 \alpha_{\mathrm{con}}^*(n)\leq n.
\]
For bideterministic input, known cycle-rank algorithms give a polynomial-time
\[
 O\bigl((\log n)^{3/2}\bigr)
\]
approximation.

\section{Hard conjecture}

\begin{conjecture}[Hard II: polylogarithmic approximation]
There is a universal constant $d$ and a polynomial-time constructive
algorithm satisfying
\[
 h(C_A)
 \leq
 O\bigl((\log n)^d\bigr)\,h(L(A))
\]
for every minimal $n$-state DFA $A$.
\end{conjecture}

This is a weaker fallback conjecture if constant approximation is too strong.
Even this would be a substantial improvement over the general linear
baseline.


\begin{thebibliography}{9}

\bibitem{Eggan}
L.~C. Eggan,
\emph{Transition Graphs and the Star-Height of Regular Events},
Michigan Mathematical Journal 10(4) (1963), 385--397.

\bibitem{Hashiguchi}
K.~Hashiguchi,
\emph{Algorithms for Determining Relative Star Height and Star Height},
Information and Computation 78(2) (1988), 124--169.

\bibitem{McNaughton1967}
R.~McNaughton,
\emph{The Loop Complexity of Pure-Group Events},
Information and Control 11(1--2) (1967), 167--176.

\bibitem{McNaughton1969}
R.~McNaughton,
\emph{The Loop Complexity of Regular Events},
Information Sciences 1(3) (1969), 305--328.

\bibitem{Gruber}
H.~Gruber,
\emph{Digraph Complexity Measures and Applications in Formal Language
Theory},
Discrete Mathematics \& Theoretical Computer Science 14(2) (2012),
189--204.

\bibitem{GruberHolzer}
H.~Gruber and M.~Holzer,
\emph{From Finite Automata to Regular Expressions and Back---A Summary on
Descriptional Complexity},
International Journal of Foundations of Computer Science 26(8) (2015),
1009--1040.

\bibitem{Kleene}
S.~C. Kleene,
\emph{Representation of Events in Nerve Nets and Finite Automata},
in \emph{Automata Studies}, Annals of Mathematics Studies 34,
Princeton University Press, 1956, 3--41.

\end{thebibliography}

\end{document}
